\begin{abstract}

This thesis explores the applicability of Kubernetes as a platform for implementing business logic using Custom Resource Definitions (CRDs) and the Operator pattern. While Kubernetes has become the de facto standard for container orchestration, its potential as a business logic platform remains understudied.

Two functionally equivalent billing systems were designed and implemented: (1) a traditional Go backend with REST API and PostgreSQL database, and (2) a Kubernetes Operator with CRDs and reconcile-based control loop. Both systems implement identical business logic (billing plan management, subscription lifecycle, recurring billing, expiration handling) but differ fundamentally in architecture, state management, and execution model.

Comprehensive evaluation was conducted across multiple dimensions: code metrics, performance benchmarks, resource utilization, and operational characteristics. The results demonstrate that the traditional backend achieves 10--20$\times$ lower latency (p50: 10ms vs 100ms for subscription creation) and 5$\times$ higher throughput (~1,000 req/s vs ~200 req/s). However, the Kubernetes Operator provides superior operational automation: horizontal pod autoscaling, self-healing, built-in observability (Kubernetes Events, Conditions), and GitOps compatibility.

The research contributes empirical evidence to the cloud-native architecture discourse, moving beyond anecdotal claims to data-driven decision making. A decision framework is provided for practitioners: the traditional approach is recommended for performance-critical, portable applications with simple deployment requirements; the Operator approach is suitable for organizations already invested in Kubernetes ecosystem that prioritize operational automation over raw performance.

Both implementations are released as open source and are presented as reference rather than production code; outstanding limitations (no payment-gateway integration, currency-agnostic revenue metric, single-replica worker, asymmetric lifecycle vocabularies between the two systems) are catalogued in the conclusion. The findings suggest that Kubernetes can serve as a viable business-logic platform when business workflows are naturally expressible as resource lifecycles, but that the latency-automation trade-off should be evaluated explicitly against organisational context, team expertise, and the specific shape of the workload, rather than treated as a technological default.

\end{abstract}

\begin{abstract}[russian]

Данная диссертация исследует применимость Kubernetes в качестве платформы для реализации бизнес-логики с использованием Custom Resource Definitions (CRD) и паттерна Operator. Несмотря на то, что Kubernetes стал стандартом де-факто для оркестрации контейнеров, его потенциал как платформы для бизнес-логики остаётся недостаточно изученным.

Были спроектированы и реализованы две функционально эквивалентные биллинг-системы: (1) традиционный Go-бэкенд с REST API и базой данных PostgreSQL, (2) Kubernetes Operator с CRD и циклом reconcile. Обе системы реализуют идентичную бизнес-логику (управление тарифными планами, жизненный цикл подписок, периодический биллинг, обработка истечения срока), но фундаментально различаются архитектурой, управлением состоянием и моделью выполнения.

Проведено комплексное сравнение по нескольким измерениям: метрики кода, бенчмарки производительности, использование ресурсов и операционные характеристики. Результаты показывают, что традиционный бэкенд обеспечивает в 10--20 раз меньшую задержку (p50: 10 мс против 100 мс для создания подписки) и в 5 раз более высокую пропускную способность (~1000 запросов/с против ~200 запросов/с). Однако Kubernetes Operator предоставляет превосходную операционную автоматизацию: автоматическое масштабирование (HPA), самовосстановление, встроенную наблюдаемость (Kubernetes Events, Conditions) и совместимость с GitOps.

Исследование вносит эмпирические данные в дискуссию об облачных нативных архитектурах, переходя от анекдотических утверждений к обоснованным решениям. Предоставлена рамка принятия решений для практиков: традиционный подход рекомендуется для критичных к производительности приложений с простыми требованиями к развёртыванию; подход с Operator подходит для организаций, уже инвестирующих в экосистему Kubernetes и приоритизирующих операционную автоматизацию над производительностью.

Обе реализации выпущены как open-source в качестве референсного, а не production-кода; остающиеся ограничения (отсутствие интеграции с платёжным шлюзом, отсутствие меток валюты в метрике выручки, отсутствие leader-election в фоновом воркере, асимметрия словарей состояний между двумя системами) систематизированы в заключении. Выводы показывают, что Kubernetes может служить жизнеспособной платформой для бизнес-логики, когда бизнес-процессы естественно выражаются через жизненный цикл ресурсов, но компромисс между задержкой и автоматизацией следует оценивать применительно к организационному контексту, экспертизе команды и конкретной форме нагрузки, а не принимать как технологическую данность.

\end{abstract}
